\chapter{Literature Review and Theoretical Background}\label{ch:literature}

% Generated by docs/research/replication/audit.py; do not edit.
\newcommand{\OIAcPools}{12}
\newcommand{\OIAcRuns}{250}
\newcommand{\OIAcSweepMax}{3.8}
\newcommand{\OIAcSweepMin}{1.2}
\newcommand{\OIDwlExploitRuns}{99}
\newcommand{\OIDwlVehicles}{5{,}860}
\newcommand{\OIFactCount}{77}
\newcommand{\OIFloorThree}{0.25}
\newcommand{\OISelectionGap}{2.50}
\newcommand{\OIVslBiasMax}{0.0064}
\newcommand{\OIVslBiasMean}{0.0026}
\newcommand{\OIVslChangesMax}{92}
\newcommand{\OIVslChangesMin}{86}
\newcommand{\OIVslGain}{5.46}
\newcommand{\OIVslJumpMax}{70}
\newcommand{\OIVslMaxSeedA}{3.61}
\newcommand{\OIVslMaxSeedB}{0.62}
\newcommand{\OIVslMaxSeedC}{-8.88}
\newcommand{\OIVslPooledMax}{-0.56}
\newcommand{\OIVslSeedA}{6.69}
\newcommand{\OIVslSeedB}{5.37}
\newcommand{\OIVslSeedC}{4.27}
\newcommand{\OIVslShareFortyMax}{44}
\newcommand{\OIVslShareFortyMin}{43}


\todo{Sections on freeway control, deep reinforcement learning in traffic and evaluation practice, from thesis outline v1.}

\section[Open implementations]{Open implementations: design choices and the evidence they support}\label{sec:open-impl}

A review of papers compares what authors report. It cannot show which traffic draws produced a number, which run was selected, or which part of the road a travel time covers. Public code and stored results can. This section therefore reads three open implementations of learning-based motorway speed control as primary sources, alongside their publications where those exist, and asks three questions: where each places learning in the control loop, what its evidence can establish, and what it would take to run it again.

The three are a purposive sample, chosen because they are public, built on SUMO and closest to the actuators studied here. They are not the result of a systematic search and no statement below generalises beyond them.

\begin{description}[leftmargin=0pt,style=nextline]
\item[ACMPC] An MPC of variable speed limits and ramp metering in which tabular Q-learning selects one model parameter online \parencite{ACMPCRepo}. A search found no associated publication, so the repository is the only source.
\item[vsl2] A two-level soft actor-critic controller on the I-80 at Emeryville \parencite{Vsl2Repo}, derived from an earlier public code base \parencite{LuSacVsl2022}. A search found no associated publication.
\item[DWL-ST-VSL] Distributed W-learning agents that adapt both the speed limit and the position of its zone \parencite{KusicRepo}, published in \textcite{Kusic2020}, \textcite{Kusic2021ITSC} and \textcite{Kusic2021Math}\verify.
\end{description}

\subsection{Method of the review}\label{sec:open-impl-method}

Each repository was fixed at one commit. A script reads the files at those commits as text, checks every statement used in this section against a specific file, and recomputes each quoted number from stored results; if any statement stops holding, it stops and writes nothing \ledger{docs/research/replication/audit.py}. The numbers printed below are generated by that script, which currently checks \OIFactCount{} statements.

Three limits apply. Nothing has been executed: this section reports what the artefacts contain, not how the controllers perform. A repository is a snapshot, and a missing file is not evidence that the authors lacked it; published work may rest on material that was never released. Finally, the publications of DWL-ST-VSL could not be retrieved in full for this draft, and are marked accordingly until they are read. The original authors have not yet been contacted.

\subsection{Where learning sits in the control loop}\label{sec:open-impl-design}

All four systems discussed here, the three above and the one studied in this thesis, close the same loop. With traffic state $x_t$, measurements $y_t$, command $u_t$, and exogenous randomness $\xi$ (demand realisation, driver draws),
\begin{equation}
x_{t+1}=f(x_t,u_t,\xi_t),\qquad u_t=\kappa\bigl(y_t;\theta_t\bigr).\label{eq:loop}
\end{equation}
They differ in what the learner is allowed to change: the parameter $\theta_t$ of a fixed law $\kappa$, or the law itself.

\paragraph{Learning a parameter of a model-based controller (ACMPC).}
A METANET prediction model in the tradition of \textcite{Hegyi2005}\verify{} is optimised over a horizon, with cost
\begin{equation}
J_{\mathrm{MPC}}=T\sum_{k}\Bigl(\sum_i \rho_i(k)L_i\lambda_i+\sum_o w_o(k)\Bigr)+\text{penalties on the variation of }v^{\mathrm{ctrl}}\text{ and }r,
\end{equation}
where $\rho_i$ is density, $L_i$ and $\lambda_i$ are segment length and lanes, and $w_o$ are origin queues. The speed limit enters through the equilibrium speed. As implemented, with $V$ a reference speed and $\alpha$ the compliance parameter,
\begin{equation}
\begin{gathered}
b=\min\Bigl(\tfrac{v^{\mathrm{ctrl}}}{V}(1+\alpha),\,1\Bigr),\qquad v^{\star}=\min(bV,\,v_f),\\
\rho^{\star}=\rho_{\mathrm{cr}}\bigl(1+A(1-b)\bigr),\qquad a^{\star}=a\bigl(E-(E-1)b\bigr).
\end{gathered}\label{eq:frejo}
\end{equation}
The repository attributes this form to Frejo and co-authors \todo{identify and read the source}; related forms appear in \textcite{Carlson2010}\verify. The learner is a tabular agent that, every \SI{240}{\second}, selects $\alpha$ from five values given one of six discrete states, by the update $Q(s,a)\leftarrow Q(s,a)+\eta\bigl(r+\gamma\max_{a'}Q(s',a')-Q(s,a)\bigr)$. The reward penalises total time spent, including ramp queues, and a clipped collision count. The learned object is thirty numbers, and the road never receives a command that the MPC did not produce.

\paragraph{Learning at two time scales (vsl2).}
A manager selects one of eight limits between 30 and \SI{100}{\kilo\metre\per\hour} every \SI{60}{\second}, applied to the lanes of one edge. A worker adds an offset within \SI{\pm 3}{\metre\per\second} every \SI{4}{\second}, sent as one common speed to all connected vehicles on that edge. Both are trained by soft actor-critic \parencite{Haarnoja2018}\verify{}. The manager is rewarded for travel time, with a term for smooth changes, and the worker for low speed variance. Connected vehicles are therefore commanded as a group: the worker has one output, not one per vehicle.

\paragraph{Arbitration between policies (DWL-ST-VSL).}
Each agent holds local and remote policies, each with a Q-table and a W-table over 4{,}608 states. Eight actions combine four speed limits (60 to \SI{120}{\kilo\metre\per\hour}) with two zone positions, and consecutive limits may differ by at most \SI{20}{\kilo\metre\per\hour}. W estimates how much a policy stands to lose if another policy's action is followed. The action executed is that of the policy with the largest nominated weight,
\begin{equation}
a^{\mathrm{win}}=\operatorname*{arg\,max}\bigl\{W^{\mathrm{LP}}_{\max},\;C\,W^{\mathrm{RP}}_{\max}\bigr\},\qquad C=0.75,
\end{equation}
every \SI{150}{\second}. What is learned here includes where control acts, not only how strongly.

\paragraph{A learned policy per vehicle (this thesis).}
The controllers studied here set the time headway of each automated vehicle individually, through one neural policy shared by the fleet \todo{cross-reference to the plant and results chapters}. The learned object is the largest of the four, and the command is the most finely resolved.

\medskip
\Cref{tab:open-design} sets the four side by side. The ordering suggests a trade-off that the rest of this thesis returns to: the smaller the learned object, the less a learning failure can cost and the less learning can gain. This can be stated exactly for the parameter case. Let $J(\pi,\xi)$ be the realised cost of a parameter schedule $\pi$ under randomness $\xi$, and let $\alpha$ range over constant parameters. The adaptation gain
\begin{equation}
G=\frac{\min_{\alpha}\mathbb{E}_{\xi}\bigl[J(\alpha,\xi)\bigr]-\mathbb{E}_{\xi}\bigl[\min_{\pi}J(\pi,\xi)\bigr]}{\min_{\alpha}\mathbb{E}_{\xi}\bigl[J(\alpha,\xi)\bigr]}\label{eq:adaptation-gain}
\end{equation}
bounds from above the relative improvement of any causal scheduler over the best constant parameter, because the inner minimum sees $\xi$ in full. ACMPC's stored open-loop sweep gives one indication of the scale involved: across eight constant values of $\alpha$, total time spent varies by \OIAcSweepMin{} to \OIAcSweepMax{}\,\% depending on the scenario. A spread over constants is not $G$, and the macroscopic model used for the sweep contains none of the model error that a microscopic plant would add, so this is a statement about sensitivity and not about the gain available in closed loop.

\begin{table}[tbp]
\centering\footnotesize\setlength{\tabcolsep}{3pt}
\caption{Where learning sits in four motorway control systems. ACMPC, vsl2 and DWL-ST-VSL are as found in the repositories at the pinned commits.}\label{tab:open-design}
\begin{tabular}{@{}>{\raggedright\arraybackslash}p{24mm}*{4}{>{\raggedright\arraybackslash}p{29mm}}@{}}
\toprule
& ACMPC & vsl2 & DWL-ST-VSL & This thesis \\
\midrule
Command on the road & Section speed limits and ramp metering, from the MPC & Limit on one edge; one common speed for connected vehicles on it & Limits per section, with variable zone position & Individual time headway per automated vehicle \\
\addlinespace
What is learned & One MPC parameter $\alpha$ & Limit level and a common offset & Limit level and zone position & The full per-vehicle policy \\
\addlinespace
Learned object & Table, $6\times5$ & Two neural policies, soft actor-critic & Tables, $4{,}608\times8$ per policy & One shared neural policy, PPO \\
\addlinespace
Decision period & \SI{240}{\second} & \SI{60}{\second} and \SI{4}{\second} & \SI{150}{\second} & \todo{from plant chapter} \\
\addlinespace
Road & Six segments and one on-ramp; called A3 in the documentation & I-80 Emeryville geometry; seven demand intervals & Described by its authors as a synthetic motorway & Synthetic r2m3, uncalibrated by decision \ledger{V245} \\
\bottomrule
\end{tabular}
\end{table}

\subsection{What an evaluation can establish}\label{sec:open-impl-validity}

Deep reinforcement learning results are known to depend strongly on seeds, implementation details and reporting practice \parencite{Henderson2018}\verify{}, and point estimates from few runs can mislead \parencite{Agarwal2021}\verify{}. The same concerns can be stated for a traffic controller $c$ compared with a baseline $b$ through the paired effect
\begin{equation}
\Delta=\operatorname{med}_{\xi}\bigl[J_c(\xi)-J_b(\xi)\bigr].\label{eq:effect}
\end{equation}
Reading \cref{eq:effect} literally gives five requirements for a claim about $\Delta$.

\begin{description}[leftmargin=0pt,style=nextline]
\item[R1: the randomness is sampled.] $\Delta$ is a statement about a distribution over $\xi$, so $\xi$ must vary. With $n$ paired runs that all favour the controller, the exact two-sided sign-flip test still returns $p=2^{1-n}$. Three pairs cannot go below $p=\OIFloorThree$, whatever the effect, and at least six pairs are needed to reach $p<0.05$.
\item[R2: selection is independent of the evaluation.] If $k$ runs from one distribution with standard deviation $\sigma$ are available and the lowest cost is reported, the reported value is optimistic in expectation by $c_k\sigma$, with $c_{\OIDwlExploitRuns}\approx\OISelectionGap$ for normal draws, even when the controller has no effect.
\item[R3: the cost is closed under displacement.] A vehicle's time decomposes as $T_i=T^{\mathrm{ins}}_i+T^{\mathrm{up}}_i+T^{\mathrm{sec}}_i+T^{\mathrm{down}}_i$: waiting to enter, upstream of the controlled area, within a measured section and downstream of it. A controller that holds vehicles back moves time from $T^{\mathrm{sec}}$ to $T^{\mathrm{up}}$ or $T^{\mathrm{ins}}$. A cost computed on $T^{\mathrm{sec}}$ alone improves even if the total does not.
\item[R4: the opponent is strong enough for the claim.] A gain over no control shows that control helps. It does not show that learning helps, since a simple tuned law may achieve the same.
\item[R5: the side effects are measured.] Efficiency can be bought with hard braking, close following or abrupt changes of the command, and these must be reported with stated tolerances.
\end{description}

\Cref{tab:open-evidence} reports what the three repositories contain against these requirements.

\begin{table}[tbp]
\centering\footnotesize\setlength{\tabcolsep}{3pt}
\caption{Evidence in the repositories at the pinned commits, against the requirements of \cref{sec:open-impl-validity}. Absence from a repository does not establish absence from the published study. The last column gives the corresponding provision of this thesis.}\label{tab:open-evidence}
\begin{tabular}{@{}>{\raggedright\arraybackslash}p{24mm}*{4}{>{\raggedright\arraybackslash}p{29mm}}@{}}
\toprule
& ACMPC & vsl2 & DWL-ST-VSL & This thesis \\
\midrule
R1 Traffic draws in evaluation & SUMO seed 42 for every repeat, unless the unreleased configuration overrides it & Three seeds (101--103) & One SUMO seed (28815), also used for training & Fresh paired seeds per round; 240 per cell planned for the current study \\
\addlinespace
R2 Policy selection & Not visible; \OIAcPools{} experiments committed & One stored checkpoint; no validation step in the training code & Analysis script takes the lowest of \OIDwlExploitRuns{} exploitation runs & Validation seeds only; test seeds used once \\
\addlinespace
R3 Scope of the cost & Time spent including ramp queues & Mean travel time on a \SI{200}{\metre} section downstream of the controlled edge & Vehicles in the network; insertion backlog not counted & Door-to-door delay with origin and insertion queues \\
\addlinespace
R4 Opponents in the code & MPC with constant $\alpha$; speed limits only; metering only; no control & No control & No speed limits; the journal paper also reports W-learning and proportional control\verify & Tuned rule controllers, adaptive scheduler, fitted-model MPC, no control \\
\addlinespace
R5 Side effects & Collisions, inside the reward & None in the environment & None in the main script & Collisions, close following and braking at \SI{0.1}{\second} \\
\addlinespace
Inputs to run again & SUMO network in no commit & Present; SUMO version not recorded & Present; absolute Windows paths & In the repository, tagged \\
\addlinespace
Stored results & Open-loop sweeps and training logs; evaluation data not included & Nine evaluation files and one checkpoint & None & All runs in the ledger \\
\addlinespace
Automated tests & None, stated by the authors & None & None & Unit tests and a draft plant verification record \\
\addlinespace
Licence & MIT & BSD 3-Clause & GPL-3.0 & --- \\
\bottomrule
\end{tabular}
\end{table}

\subsection{Findings per implementation}\label{sec:open-impl-findings}

\paragraph{ACMPC.}
This is the implementation from which this thesis can learn most. It runs an optimisation-based controller in closed loop with a microscopic simulator on a separate thread, so that the time a solve takes is part of the experiment. It gives learning a narrow, interpretable task on top of a controller that can already act. Its opponent for the learner is the same controller with a constant parameter, which isolates the value of adaptation; its reward counts ramp queues and collisions; and its documentation lists its own defects plainly, including the absence of tests and unseeded exploration.

Several points limit what can be concluded from the repository alone. The SUMO network was never committed, so the microscopic experiments cannot be repeated from it. Every repeat is launched with seed 42 unless the unreleased configuration enables SUMO's random seeding; if it does not, repeats at constant $\alpha$ differ only through the timing of the solver thread, and the bootstrap intervals describe that timing and not traffic variability. The configuration of the final training uses a state feature that the current commit no longer implements; it is present at an earlier commit, so that experiment can still be targeted, at that commit. The evaluation data behind the report are not included, and no calibration of the prediction model against the microscopic plant appears in the code. Whether adaptive $\alpha$ improves on constant $\alpha$ in closed loop therefore remains open from this source.

\paragraph{vsl2.}
The value of this implementation is architectural: it combines a limit for all drivers with a separate channel for connected vehicles, and it was designed for a sweep over connected-vehicle share. Its stored results allow an independent check. Recomputing from the nine stored evaluation files reproduces its summary table: at 10\,\% connected vehicles, mean travel time improves by \OIVslGain{}\,\% over no control (\OIVslSeedA{}, \OIVslSeedB{} and \OIVslSeedC{}\,\% in the three seeds), and the highest travel time changes by $\OIVslPooledMax{}$\,\% when pooled ($+\OIVslMaxSeedA{}$, $+\OIVslMaxSeedB{}$ and $\OIVslMaxSeedC{}$\,\% per seed), with positive values meaning improvement.

The same files show three things the table does not. First, the connected-vehicle layer is effectively inactive in this policy: the worker's offset averages \SI{\OIVslBiasMean}{\metre\per\second} and never exceeds \SI{\OIVslBiasMax}{\metre\per\second}, against a permitted range of \SI{\pm 3}{\metre\per\second}. The reported effect is therefore consistent with the sign layer alone. Second, the reported travel time covers \SI{200}{\metre} downstream of the controlled edge, while the sign holds \SI{40}{\kilo\metre\per\hour} for \OIVslShareFortyMin{} to \OIVslShareFortyMax{}\,\% of the run; time spent on the controlled edge is outside the reported number (R3), although it is inside the manager's training reward. Third, the limit changes \OIVslChangesMin{} to \OIVslChangesMax{} times per run at one decision per minute, with steps of up to \SI{\OIVslJumpMax}{\kilo\metre\per\hour} between consecutive decisions (R5). With three seeds and no control as the only opponent, R1 and R4 are also not met. Only the 10\,\% policy is stored, so the sweep it was built for is not available.

\paragraph{DWL-ST-VSL.}
This is the line of work that treats the position of the control zone as a decision variable, and the only one of the three for which publications were found and whose repository holds a complete set of simulation inputs. Its action set is physically constrained, with consecutive limits at most \SI{20}{\kilo\metre\per\hour} apart. Its abstract reports better performance than W-learning control and a proportional speed controller on static zones, at medium and high traffic loads \parencite{Kusic2021Math}\verify{}. The latter comparison is not in the public code, which covers the two-agent, medium-load scenario.

In the public code, training and evaluation share one SUMO seed and one route file of \OIDwlVehicles{} vehicles, so traffic randomness is not sampled (R1). The analysis script reports the lowest total time spent among \OIDwlExploitRuns{} runs without exploration and compares it with a single run without speed limits (R2). \textcite{Kusic2021Math}\verify{} may report differently; this has to be checked against the full text before anything is concluded about the published numbers. \todo{Read the three papers: replication design, comparators, reported statistics.}

\subsection{Consequences for this thesis}\label{sec:open-impl-consequences}

The review supports the design of the evaluation method of this thesis \todo{cross-reference to the method chapter}, and each requirement corresponds to a provision there: paired fresh seeds (R1), selection on validation seeds only with a single use of test seeds (R2), door-to-door delay with origin and insertion queues (R3), tuned non-learning controllers in every table (R4) and step-resolution safety measures with stated tolerances (R5). These provisions are a response to known failure modes and to errors found during this project \ledger{V244, V283, V294, V367, V386}. They do not by themselves make the results here more correct; they make them easier to check.

The review also identifies where these implementations go further than this thesis, and those points are limits of the present work:
\begin{enumerate}[nosep]
\item Physics-based MPC in closed loop. The model-based opponent used here selects among fixed actions with a fitted neural surrogate \ledger{planb.md \S3.5}. It does not test prediction from a traffic-flow model or the effect of solve latency.
\item Learning as a bounded supervisor of a working controller. A per-vehicle residual has been tried here and did not pass validation \ledger{V449}, but a learned scheduler of a tuned rule's parameters has not.
\item Coordination with ramp metering. Mainline metering was closed on this project's plant \ledger{V149}, so coordinated actuation is outside the claim.
\item A sweep over automated-vehicle share. The current study fixes one share.
\item A named road. The plant here is synthetic by decision \ledger{V245}. None of the three repositories shows a calibration either, but a named corridor gives an external reference that a synthetic plant lacks.
\end{enumerate}

\subsection{A replication programme}\label{sec:open-impl-replication}

Reading an implementation shows what it claims to do, while running it shows whether the claim holds and whether this project can operate another group's controller at all. The programme therefore climbs an evidence ladder, in which each level is a precondition of the next (\cref{tab:open-ladder}). Pass criteria are fixed before any run, in a protocol kept with the other protocols of this project \ledger{docs/research/protocols/replication\_v0.md}. To date only the two levels that require no simulation have been carried out, and \emph{no controller from these repositories has been executed by this project}.

\begin{table}[tbp]
\centering\footnotesize\setlength{\tabcolsep}{3pt}
\caption{Evidence ladder and current status. Open entries are planned, not results.}\label{tab:open-ladder}
\begin{tabular}{@{}l>{\raggedright\arraybackslash}p{48mm}*{3}{>{\raggedright\arraybackslash}p{20mm}}@{}}
\toprule
Level & Question & ACMPC & vsl2 & DWL-ST-VSL \\
\midrule
E0 Inventory & What exists, at which commit, under which licence? & Done & Done & Done \\
E1 Recomputation & Do stored results give the reported numbers, without simulation? & Partial: sweeps only & Done: table reproduced & Not possible: no stored results \\
E2 Execution & Does the pipeline run at its declared versions, and does a stored seed repeat? & Blocked: network missing & Planned & Planned \\
E3 Claim reproduction & Does the reported effect appear with fresh seeds on the same plant? & Needs a reported effect & Planned & Planned \\
E4 Common protocol & Same controller and plant under R1--R5, against tuned opponents? & Planned & Planned & Planned \\
E5 Transfer & Does the idea survive on this thesis's plant? & Optional & Optional & Optional \\
\bottomrule
\end{tabular}
\end{table}

\todo{Fill E2 to E4 as they are executed; every open cell stays open until a run in the ledger supports it.}
